% Bangla report for the heavy Initial Solution V1 experiment.
% Compile with XeLaTeX so Bengali text is shaped correctly.
\documentclass[10pt,a4paper]{article}
\usepackage{fontspec}
\usepackage{polyglossia}
\setmainlanguage{bengali}
\setotherlanguage{english}
\setmainfont[Script=Bengali]{Kohinoor Bangla}
\setsansfont[Script=Bengali]{Kohinoor Bangla}
\newfontfamily\bengalifonttt[Script=Bengali]{Kohinoor Bangla}
\usepackage[a4paper,margin=22mm,headheight=24pt]{geometry}
\usepackage{amsmath}
\usepackage{array}
\usepackage{booktabs}
\usepackage{enumitem}
\usepackage{fancyhdr}
\usepackage{longtable}
\usepackage{pdflscape}
\usepackage{tabularx}
\usepackage[table]{xcolor}
\usepackage{hyperref}

\definecolor{navy}{HTML}{17365D}
\definecolor{bluegray}{HTML}{EAF1F8}
\definecolor{lightgray}{HTML}{F4F6F8}
\definecolor{midgray}{HTML}{5B6573}
\definecolor{rulegray}{HTML}{C9D1D9}
\hypersetup{
  colorlinks=true,
  linkcolor=navy,
  urlcolor=navy,
  pdftitle={Initial Solution V1 Simulation Report},
  pdfauthor={PGBM Thesis Simulation Work}
}
\setlist[itemize]{leftmargin=*,itemsep=2pt,topsep=3pt}
\setlist[enumerate]{leftmargin=*,itemsep=4pt,topsep=3pt}
\setlength{\parindent}{0pt}
\setlength{\parskip}{6pt}
\setlength{\emergencystretch}{1em}
\pagestyle{fancy}
\fancyhf{}
\fancyhead[L]{\textcolor{midgray}{PGBM Thesis Simulation}}
\fancyhead[R]{\textcolor{midgray}{Initial Solution V1}}
\fancyfoot[C]{\textcolor{midgray}{\thepage}}
\renewcommand{\headrulewidth}{0.4pt}
\renewcommand{\headrule}{\hbox to\headwidth{\color{rulegray}\leaders\hrule height \headrulewidth\hfill}}
\newcommand{\code}[1]{\path{#1}}

\begin{document}

\begin{titlepage}
\thispagestyle{empty}
\vspace*{18mm}
{\color{navy}\rule{\textwidth}{1.5pt}}\\[10mm]
{\color{navy}\sffamily\bfseries\fontsize{25}{32}\selectfont Initial Solution V1\\Simulation Report\par}
\vspace{6mm}
{\color{midgray}\sffamily\Large Simulation-এর core feature, algorithm, experiment এবং ফলাফল\par}
\vfill
\begin{tabular}{@{}ll@{}}
\textbf{Project:} & PGBM thesis and simulation work\\
\textbf{সমাধানের ধরন:} & Bounded brute-force reference planner\\
\textbf{Environment:} & Synthetic এবং DU outdoor\\
\textbf{মোট experiment:} & ২৪০টি run\\
\textbf{Report date:} & ২৭ সেপ্টেম্বর ২০২৬
\end{tabular}
\vfill
\begin{center}
\colorbox{bluegray}{\begin{minipage}{0.88\textwidth}
এই report-এ Initial Solution V1-এর simulation কীভাবে কাজ করে, assignment এবং routing কীভাবে নির্ধারিত হয়, কোন experiment চালানো হয়েছে এবং সেই experiment থেকে কী ফলাফল পাওয়া গেছে তা সরাসরি ব্যাখ্যা করা হয়েছে।
\end{minipage}}
\end{center}
\vspace{12mm}
{\color{navy}\rule{\textwidth}{1.5pt}}
\end{titlepage}

\tableofcontents
\clearpage

\section{Simulation-এর মূল বৈশিষ্ট্য}

আমাদের simulation একটি static একবারের assignment নয়। এটি একটি seeded, event-driven simulation, যেখানে ১২০ মিনিটের মধ্যে task আসে, queue-তে জমা হয়, available UAV-এর জন্য plan তৈরি হয়, route execute হয় এবং UAV আবার base-এ ফিরে আসে।

Simulation-এর প্রধান বৈশিষ্ট্যগুলো হলো:

\begin{itemize}
\item \textbf{১২০ মিনিটের operation window:} পুরো simulation ০ থেকে ১২০ মিনিট পর্যন্ত চলে।
\item \textbf{Dynamic task arrival:} সব task শুরুতেই একসাথে available থাকে না। নির্দিষ্ট সময় অনুযায়ী task detect হয়।
\item \textbf{তিনটি arrival phase:} ০--৪০ মিনিটে high, ৪০--৮০ মিনিটে medium এবং ৮০--১২০ মিনিটে low arrival phase থাকে।
\item \textbf{Task queue:} সব UAV ব্যস্ত থাকলে নতুন task queue-তে অপেক্ষা করে।
\item \textbf{দুইটি environment:} Synthetic এবং DU outdoor environment-এ একই ধরনের experiment চালানো হয়েছে।
\item \textbf{Individual task model:} প্রতিটি task আলাদা demand হিসেবে model করা হয়েছে। অতিরিক্ত household বা building-level demand model ব্যবহার করা হয়নি।
\item \textbf{Task এবং parcel আলাদা:} প্রতিটি task-এ food, water এবং medical item-এর subset থাকতে পারে। তাই একটি task-এ একাধিক parcel থাকতে পারে।
\item \textbf{Fixed drop-off time:} একটি parcel drop করতে ৩০ সেকেন্ড বা ০.৫ মিনিট ধরা হয়েছে। একটি task-এর drop-off time শুধু parcel সংখ্যার উপর নির্ভর করে। Severity drop-off time পরিবর্তন করে না।
\item \textbf{Drop-off visualization:} UAV base থেকে task-এর vertical drop-off waypoint-এ যায়, সেখানে নির্ধারিত service hold করে, তারপর পরবর্তী waypoint বা base-এ ফেরে। এই arrival, hold এবং return event log-এ record করা হয়।
\item \textbf{Seeded reproducibility:} একই seed এবং একই configuration ব্যবহার করলে একই scene এবং task set পুনরায় তৈরি হয়।
\item \textbf{No recourse:} কোনো UAV active mission-এ থাকা অবস্থায় নতুন task-এর জন্য route পরিবর্তন করে না।
\end{itemize}

\subsection{Environment এবং input provenance}

V1-এর দুইটি environment একই ধরনের নয়, তাই তাদের result একই অর্থে পড়া যাবে না।
Synthetic environment-এ building layout, damage geometry এবং collision obstacle
procedurally তৈরি হয়। DU outdoor environment-এ local OSM-derived DU Science
Complex building footprint এবং road metadata ব্যবহার করা হয়; তার উপর earthquake
damage এবং blocked-volume layer simulation-এর জন্য derived হিসেবে তৈরি করা হয়।
অতএব DU outdoor case-কে actual post-earthquake survey বা measured field
reconstruction বলা যাবে না। এটি map-informed এবং derived-damage environment।

DU-এর road geometry source context হিসেবে সংরক্ষিত থাকলেও airborne UAV route
check-এ road-কে active collision obstacle করা হয়নি। V1-এর route difficulty মূলত
building এবং derived obstacle geometry থেকে এসেছে। এই provenance জানা জরুরি,
কারণ DU outdoor-এর কম service rate বা বেশি route time-কে universal real-world
performance হিসেবে interpret করা যাবে না।

\subsection{Task generation-এর বিস্তারিত নিয়ম}

V1-এ task কোনো household বা পুরো building-এর aggregate demand নয়; প্রতিটি task
একটি individual survivor-assistance delivery request। Task location minor, major
বা destroyed damage state-যুক্ত target-building footprint-এর ভিতরে sample করা
হয়। UAV সরাসরি survivor position-এ নামে না; survivor position-এর একই
horizontal location-এ ২--৪ মিটার উপরের drop-off waypoint ব্যবহার করে।

\begin{table}[htbp]
\centering
\caption{V1 task input generation-এর প্রধান নিয়ম।}
\label{tab:task-generation-bn}
\scriptsize
\rowcolors{2}{lightgray}{white}
\begin{tabular}{@{}p{0.28\textwidth}p{0.61\textwidth}@{}}
\toprule
বিষয় & V1-এ ব্যবহৃত নিয়ম \\
\midrule
Relief item & Food, water এবং medical; প্রতিটি task-এ অন্তত একটি item থাকে \\
Item quantity & প্রতিটি item-এর quantity ০ অথবা ১; তাই একটি task-এ সর্বোচ্চ ৩টি parcel থাকে \\
Unit weight & Food ০.২৫ kg, water ০.৫০ kg, medical ০.৫০ kg \\
Severity & ০.৩৫ থেকে ১.০০ range-এর মধ্যে generated operational priority score \\
Detection time & High, medium এবং low phase-এর quota অনুযায়ী generated detection time \\
Drop-off point & Survivor position-এর উপরে ২--৪ m vertical separation-এ UAV-accessible point \\
\bottomrule
\end{tabular}
\end{table}

এই নিয়মের কারণে task count এবং parcel count এক নয়। একই সঙ্গে ২ kg UAV payload
capacity একাধিক individual task-এর parcel একসাথে carry করার সময় পরীক্ষা হয়;
একটি একক task-এর demand দিয়ে পুরো payload capacity সাধারণত পূর্ণ হয় না।

V1-এর উদ্দেশ্য হলো একটি পরিষ্কার এবং পরিমাপযোগ্য reference solution তৈরি করা। তাই এটি final mathematical PGBM optimizer হিসেবে দাবি করা হচ্ছে না।

\section{V1 algorithm কীভাবে assignment এবং routing ঠিক করে}

যখন কোনো UAV available হয় এবং queue-তে pending task থাকে, তখন planner একটি নতুন dispatch decision নেয়। সেই decision-এ assignment এবং routing আলাদা আলাদা ধাপে নির্ধারিত না হয়ে একসাথে বিবেচিত হয়।

প্রতিটি dispatch opportunity-তে algorithm-এর flow হলো:

\begin{enumerate}
\item Queue থেকে eligible pending task নেওয়া হয়।
\item কোন task কোন UAV serve করবে তার সম্ভাব্য combination তৈরি করা হয়।
\item প্রতিটি UAV-এর জন্য task order বা visit order তৈরি করা হয়।
\item প্রতিটি candidate-এর জন্য 3D route তৈরি এবং route feasibility check করা হয়।
\item Payload, parcel count, battery reserve, drop-off time, horizon এবং base-এ return করা সম্ভব কি না পরীক্ষা করা হয়।
\item Feasible candidate-গুলোর service value হিসাব করা হয়।
\item সর্বোচ্চ মোট service value পাওয়া candidate plan নির্বাচন করা হয়।
\end{enumerate}

Candidate plan তুলনা করার সময় V1-এর priority rule-টি lexicographic। প্রথমে
যে plan-এর total time-dependent service value বেশি সেটি অগ্রাধিকার পায়। দুইটি
plan-এর value সমান হলে বেশি task serve করা plan নেওয়া হয়। তাও সমান হলে কম
mission energy এবং শেষে কম route distance-কে tie-breaker হিসেবে ব্যবহার করা
হয়। তাই V1-এর primary objective service value; task count, energy এবং distance
মূল objective-এর বিকল্প নয়, বরং সমমানের candidate আলাদা করার secondary rule।

\subsection{Route কীভাবে তৈরি হয়}

প্রতিটি candidate plan-এ একটি UAV-এর route এই ধারায় তৈরি হয়:

\begin{equation*}
\text{Base} \;\rightarrow\; \text{Task 1 waypoint} \;\rightarrow\; \text{Task 2 waypoint} \;\rightarrow\; \cdots \;\rightarrow\; \text{Base}.
\end{equation*}

Route তৈরির বিষয়টি সহজভাবে বোঝার জন্য প্রতিটি mission-কে Base থেকে task waypoint-এ যাওয়া, parcel drop করা এবং শেষে Base-এ ফেরার একটি 3D পথ হিসেবে ভাবতে হবে। Route তৈরির ধাপগুলো হলো:

\begin{enumerate}
\item Candidate plan যে task order দিয়েছে, সেই order অনুযায়ী Base, task waypoint এবং পরবর্তী waypoint-এর মধ্যে আলাদা route leg তৈরি করা হয়।
\item প্রতিটি leg-এর জন্য algorithm প্রথমে একটি preferred cruise altitude পরীক্ষা করে। Cruise altitude হলো UAV-এর horizontal flight-এর জন্য ব্যবহৃত মাঝ-আকাশের উচ্চতা; এটি task waypoint-এর ২--৪ মিটার drop-off height নয়। V1-এ preferred cruise altitude ১২ মিটার, তবে endpoint বা obstacle-এর কারণে এটি ব্যবহার করা না গেলে higher fallback altitude-ও পরীক্ষা করা হয়।
\item ওই altitude-এ start point থেকে target point পর্যন্ত সরাসরি horizontal line কোনো building বা polygon obstacle-এর ভিতর দিয়ে যাচ্ছে কি না দেখা হয়। এটিই cruise-altitude test-এর অর্থ। পথটি clear হলে direct segment নেওয়া হয়।
\item Direct segment blocked হলে ৫ মিটার grid এবং আট দিকের প্রতিবেশী ব্যবহার করে A* search চালানো হয়, যাতে obstacle ঘুরে একটি collision-free horizontal path পাওয়া যায়।
\item Start point থেকে cruise altitude-এ ওঠা এবং target point-এ নামার vertical segment-ও obstacle-এর সঙ্গে intersect করছে কি না পরীক্ষা করা হয়।
\item সব leg জোড়া লাগিয়ে Base থেকে task waypoint-গুলো ঘুরে আবার Base-এ ফেরার সম্পূর্ণ 3D polyline তৈরি করা হয়। প্রতিটি segment-এর distance যোগ করে route distance, travel time এবং energy হিসাব করা হয়।
\end{enumerate}

তাই route শুধু Base এবং task-এর মধ্যে straight-line distance নয়; এটি obstacle-aware, altitude-aware এবং return-to-base route। কোনো candidate-এর route তৈরি করা না গেলে সেই candidate বাদ দেওয়া হয়। একই environment, দুই endpoint এবং একই routing configuration-এর leg পুনরায় লাগলে V1-এর basic in-memory route cache আগের route result পুনরায় ব্যবহার করতে পারে। এই cache route construction-এর পুনরাবৃত্তি কমায়, তবে V1-এ cache hit rate আলাদা metric হিসেবে record করা হয়নি।

এটি bounded brute-force approach। প্রতি dispatch-এ queue-র সব task নিয়ে অসীম search করা হয় না। Service value অনুযায়ী সর্বোচ্চ ৮টি pending task candidate search-এ নেওয়া হয়; বাকি task পরবর্তী dispatch-এর জন্য queue-তে থাকে। এরপর assignment এবং task order-এর সর্বোচ্চ ৫০,০০০টি candidate plan evaluate করা হয়। এই bound-এর utility হলো runtime নিয়ন্ত্রণে রাখা; কারণ UAV, task এবং order বাড়লে combination দ্রুত বেড়ে যায়। Effect হলো, ৮টির বেশি pending task থাকলে lower-ranked task ওই dispatch-এ বিবেচিত নাও হতে পারে, আর ৫০,০০০ limit-এ পৌঁছালে V1 global optimum দাবি করতে পারে না। Search limit-এ পৌঁছানো dispatch আলাদাভাবে report করা হয়েছে।

Task-এর completion time যত দেরি হয়, service value তত কমে। V1-এ ব্যবহৃত value function হলো:

\begin{equation}
v_i(C_i) = s_i \exp\!\left[-\lambda\left(C_i-t_i\right)\right].
\end{equation}

এখানে $s_i$ হলো task-এর severity, $t_i$ হলো task detect হওয়ার সময় এবং $C_i$ হলো task completion time। তাই algorithm শুধু বেশি task serve করার চেষ্টা করে না, বরং গুরুত্বপূর্ণ task যত দ্রুত সম্ভব complete করার চেষ্টা করে।

এই decay rate-টি arbitrary runtime parameter হিসেবে নেওয়া হয়নি। আগের service-value
calibration experiment-এ common exponential function severity ordering বজায় রেখে
reference policy-এর সঙ্গে সবচেয়ে consistent হওয়ায় সেটি ব্যবহার করা হয়েছে। V1-এ
$\lambda=0.011997\ \text{min}^{-1}$; এর equivalent half-life প্রায় ৫৭.৮ মিনিট,
অর্থাৎ একই severity-এর task-এর operational service value প্রায় ৫৭.৮ মিনিট পরে
অর্ধেক হয়। এই value-টি priority বা service benefit-এর model, medical survival
probability নয়।

\section{Simulation কীভাবে চলে}

একটি run-এর সময় simulation এই ধারায় এগোয়:

\begin{enumerate}
\item নির্দিষ্ট seed দিয়ে scene, base, candidate location এবং task location তৈরি হয়।
\item সময় এগোলে high, medium এবং low phase অনুযায়ী নতুন task detect হয়।
\item কোনো UAV available থাকলে planner queue থেকে task নিয়ে নতুন assignment এবং route তৈরি করে।
\item UAV base থেকে উড়ে task-এর vertical drop-off waypoint-এ যায়।
\item নির্ধারিত parcel drop-off time শেষ হলে UAV পরবর্তী waypoint বা return route অনুসরণ করে base-এ ফিরে আসে।
\item Base-এ ফিরে UAV-এর turnaround বা resupply/recharge সময় ধরা হয় ৫ মিনিট।
\item UAV আবার available হলে queue থেকে pending task নিয়ে পরবর্তী dispatch decision হয়।
\end{enumerate}

৫ মিনিট turnaround বলতে বোঝানো হচ্ছে, UAV একটি mission শেষ করে base-এ ফেরার পর পরবর্তী mission-এর জন্য প্রস্তুত হতে ৫ মিনিট সময় নেয়। এটি বর্তমান V1-এর operational assumption, কোনো field-calibrated physical measurement নয়।

প্রতিটি run-এর শুরুতে সব UAV একই physical Base-এ available অবস্থায় থাকে, full
energy capacity এবং empty onboard inventory নিয়ে। একটি mission চলাকালীন UAV
busy থাকে; Base-এ ফিরে ৫ মিনিট turnaround শেষ হলে তাকে আবার available ধরা হয়।
তবে V1-এর event simulator battery percentage বা item inventory প্রতি event-এ
স্থায়ীভাবে update করে না; পরবর্তী mission-এর feasibility aggregate mission
energy এবং configured reserve check দিয়ে নির্ধারিত হয়। তাই এই lifecycle-কে
stateful battery-and-inventory simulation হিসেবে interpret করা যাবে না।

একবার UAV-এর mission শুরু হলে নতুন task এলেও তার route পরিবর্তন হয় না। নতুন task queue-তে অপেক্ষা করে এবং পরবর্তী dispatch opportunity-এর জন্য থাকে।

\subsection{Drop-off time এবং simulation visualization}

V1-এ একটি task এবং একটি parcel একই জিনিস নয়। Task হলো একটি survivor delivery request, আর parcel হলো সেই request-এর food, water বা medical item-এর একটি unit। তাই task-এর demand-এ যতগুলো item আছে, drop-off time ততগুলো fixed parcel slot দিয়ে নির্ধারিত হয়:

\begin{equation*}
T_i^{\mathrm{drop}} = n_i^{\mathrm{parcel}} \times \tau_{\mathrm{parcel}}, \qquad \tau_{\mathrm{parcel}} = 0.5\ \text{minute}.
\end{equation*}

অর্থাৎ ১, ২ এবং ৩টি parcel-এর জন্য drop-off time যথাক্রমে ০.৫, ১.০ এবং ১.৫ মিনিট। এখানে severity priority এবং completion service value-তে ব্যবহৃত হয়, কিন্তু parcel নামানোর সময় বাড়ায় না। Drop-off waypoint survivor position-এর সরাসরি উপরে ২ থেকে ৪ মিটার vertical separation-এ থাকে। Route visualization-এ UAV প্রথমে ওই waypoint-এ পৌঁছায়, তারপর এই duration-এর জন্য service hold করে এবং শেষে route-এর পরবর্তী point-এ যায়। Base-এ ফেরার পরের ৫ মিনিট turnaround এই drop-off time-এর অংশ নয়।

\section{Experiment design}

Heavy V1 experiment-এ তিনটি বিষয় আলাদাভাবে পরিবর্তন করা হয়েছে: environment, task load এবং UAV সংখ্যা। প্রতিটি configuration-এর জন্য ১০টি seed ব্যবহার করা হয়েছে।

\begin{table}[htbp]
\centering
\caption{V1 experiment-এর মূল configuration।}
\label{tab:protocol-bn}
\rowcolors{2}{lightgray}{white}
\begin{tabularx}{0.95\textwidth}{@{}>{\bfseries}p{0.29\textwidth}X@{}}
\toprule
বিষয় & Configuration \\
\midrule
Simulation horizon & ১২০ মিনিট \\
Arrival phase & ০--৪০ high, ৪০--৮০ medium, ৮০--১২০ low \\
Phase ratio & ৩:২:১ \\
Task load & ৩০, ৬০ এবং ৯০ task \\
UAV সংখ্যা & ৩, ৫, ৮ এবং ১০ \\
Seed & ১০১ থেকে ১১০, মোট ১০টি seed \\
Payload & প্রতি UAV ২ kg এবং সর্বোচ্চ ৪ parcel \\
Base turnaround & ৫ মিনিট \\
Search bound & সর্বোচ্চ ৮টি pending task এবং ৫০,০০০ candidate \\
Recourse & V1-এ disabled \\
\bottomrule
\end{tabularx}
\end{table}

সর্বমোট run সংখ্যা:

\begin{equation*}
2\;\text{environment} \times 3\;\text{task load} \times 4\;\text{UAV count} \times 10\;\text{seed} = 240\;\text{runs}.
\end{equation*}

Phase quota প্রতিটি run-এ যাচাই করা হয়েছে। ৩০ task-এর জন্য quota ১৫/১০/৫, ৬০ task-এর জন্য ৩০/২০/১০ এবং ৯০ task-এর জন্য ৪৫/৩০/১৫।

\section{যে ফলাফলগুলো মাপা হয়েছে}

Simulation শুধু কতগুলো task serve হয়েছে তা মাপেনি। একটি task-এর মধ্যে একাধিক parcel থাকতে পারে বলে task-level এবং parcel-level দুটো result আলাদাভাবে record করা হয়েছে। Algorithm-এর service এবং computational behaviour বোঝার জন্য নিচের metric-গুলো record করা হয়েছে।

\begin{table}[htbp]
\centering
\caption{ব্যবহৃত প্রধান metric এবং তাদের অর্থ।}
\label{tab:metrics-bn}
\rowcolors{2}{lightgray}{white}
\begin{tabularx}{0.95\textwidth}{@{}>{\bfseries}p{0.29\textwidth}X@{}}
\toprule
Metric & কী বোঝায় \\
\midrule
Objective value & Completion time অনুযায়ী মোট service value \\
Served task & কতগুলো task সফলভাবে complete হয়েছে \\
Service rate & মোট task-এর মধ্যে কত শতাংশ serve হয়েছে \\
Requested parcel & সব task মিলিয়ে মোট প্রয়োজনীয় parcel সংখ্যা \\
Dropped parcel & সফলভাবে complete হওয়া task-গুলোর parcel সংখ্যা \\
Parcel delivery rate & Requested parcel-এর মধ্যে dropped parcel-এর শতাংশ \\
Deferred task & Simulation শেষে queue বা pending অবস্থায় থাকা task \\
Peak queue & Simulation চলাকালীন সর্বোচ্চ queue size \\
Completion delay & Task detect হওয়ার পর complete হতে গড় সময় \\
Route distance & UAV-গুলোর মোট travel distance \\
Total travel time & পুরো run-এ সব UAV mission-এর cumulative flight time; এটি ১২০ মিনিটের wall-clock horizon নয় \\
Total energy & পুরো ১২০ মিনিটে সব dispatched mission-এর cumulative energy \\
Payload utilization & একটি dispatch-এ পুরো fleet-এর parcel-count capacity-এর সর্বোচ্চ ব্যবহারের ratio; এটি payload-mass utilization নয় \\
Safe-return rate & যেসব dispatched mission route ধরে Base-এ ফিরেছে, তার proportion \\
Dispatch count & পুরো run-এ নতুন mission plan তৈরি হওয়ার সংখ্যা \\
Runtime & পুরো run শেষ হতে computer-এর সময় \\
Planner time & Assignment এবং route search-এ planner-এর সময় \\
Candidate count & কতগুলো possible plan evaluate করা হয়েছে \\
Route share & Planner time-এর কত অংশ route evaluation-এ গেছে \\
Search truncation & Search bound-এ পৌঁছানো dispatch-এর সংখ্যা \\
\bottomrule
\end{tabularx}
\end{table}

এখানে \textbf{গড় delay} বলতে task detect হওয়ার সময় থেকে task service complete হওয়ার সময় পর্যন্ত গড় simulation time বোঝানো হয়েছে।

\textbf{Run runtime} বলতে একটি সম্পূর্ণ experiment run শেষ করতে computer-এর wall-clock time বোঝানো হয়েছে। এর মধ্যে task generation, event simulation, dispatch decision, route search, execution এবং metrics record করার সময় অন্তর্ভুক্ত। অন্যদিকে \textbf{planner time} শুধু dispatch-এর assignment এবং candidate evaluation অংশের সময়।

\begin{table}[htbp]
\centering
\caption{Execution-level metric-এর observed summary।}
\label{tab:execution-summary-bn}
\scriptsize
\rowcolors{2}{lightgray}{white}
\begin{tabular}{@{}lrrrr@{}}
\toprule
Environment & Safe-return & Dispatch/run & Cumulative flight time/run (min) & Payload utilization \\
\midrule
DU outdoor & 100.0\% & 23.63 & 592.82 & 18.67\% \\
Synthetic & 100.0\% & 29.47 & 504.14 & 18.69\% \\
\bottomrule
\end{tabular}
\end{table}

এই table-এ ২৪০টি heavy-matrix run-এর environment-wise average দেওয়া হয়েছে।
দুই environment-এই safe-return rate ১০০ শতাংশ, কারণ planner route তৈরি করার
সময় return-to-base এবং energy-reserve feasibility না মিললে candidate বাদ দেয়;
এটি physical UAV reliability বা field safety probability নয়। Synthetic-এ
dispatch বেশি হলেও route ছোট হওয়ায় cumulative flight time DU outdoor-এর চেয়ে
কম। Cumulative flight time সব UAV ও সব mission যোগ করে হিসাব করা, তাই এটি
১২০ মিনিটের simulation clock-এর সঙ্গে সরাসরি তুলনীয় নয়।

Payload utilization-এর ১৮.৭ শতাংশ value-কে mass-এর ১৮.৭ শতাংশ ব্যবহার হিসেবে
পড়া যাবে না। বর্তমান V1 metric fleet-এর মোট parcel-count capacity-র বিপরীতে
একটি run-এর সর্বোচ্চ dispatch-level parcel usage মাপে। Payload mass, per-UAV
load এবং mission-wise utilization আলাদা করে জানাতে হলে V2-তে আরও precise
utilization metric প্রয়োজন হবে।

\subsection{Energy model এবং energy flow}

V1-এ energy মূলত একটি mission-এর feasibility constraint হিসেবে ব্যবহৃত হয়েছে। অর্থাৎ planner কোনো candidate route নির্বাচন করার আগে পরীক্ষা করে যে UAV mission শেষ করে reserve energy রেখে Base-এ ফিরতে পারবে কি না। এই model-এ প্রতিটি UAV-এর energy capacity ১০০০ J এবং reserve energy ২০০ J ধরা হয়েছে।

একটি mission-এর energy হিসাব করা হয়:

\begin{equation}
E_{\mathrm{mission}} = \sum_{\ell}\left(d_{\ell}e_{\mathrm{travel}} + d_{\ell}m_{\ell}e_{\mathrm{payload}} + a_{\ell}e_{\mathrm{ascent}}\right) + \sum_i \tau_i e_{\mathrm{service}}.
\end{equation}

এখানে $d_{\ell}$ হলো route leg-এর distance, $m_{\ell}$ হলো ঐ leg-এ বহন করা payload mass, $a_{\ell}$ হলো positive ascent distance এবং $\tau_i$ হলো task service বা parcel drop-off duration। V1 configuration-এ $e_{\mathrm{travel}}=1.0$ J/m, $e_{\mathrm{payload}}=0.05$ J/(kg m), $e_{\mathrm{ascent}}=0.20$ J/m এবং $e_{\mathrm{service}}=0.50$ J/min।

Energy flow-টি এভাবে কাজ করে:

\begin{enumerate}
\item UAV mission শুরুতে payload এবং available energy নিয়ে Base থেকে উড়ে।
\item Take-off, ascent এবং route-এর প্রতিটি leg-এর জন্য energy খরচ হয়। Full payload বহন করার সময় payload-related energy বেশি হয়।
\item একটি task-এর parcel drop করার পর carried mass কমে যায়; তাই পরবর্তী leg-এ payload energy কমে।
\item Drop-off hold এবং service duration-এর জন্য service energy যোগ হয়।
\item Base-এ ফেরার route-এর energy যোগ করে মোট mission energy বের করা হয়।
\item Candidate feasible ধরা হয় যদি $E_{\mathrm{mission}}+200 \leq 1000$ হয়। শর্তটি পূরণ না হলে candidate বাদ দেওয়া হয়।
\item Base-এ ফেরার পর ৫ মিনিটের turnaround/resupply শেষে UAV-কে পরবর্তী mission-এর জন্য available ধরা হয়।
\end{enumerate}

বর্তমান V1-এ battery percentage প্রতি event-এ আলাদা করে update করা হয় না; energy aggregate mission level-এ হিসাব করা হয়। তাই নিচের \textit{total energy} হলো পুরো ১২০ মিনিটে সব UAV এবং সব dispatched mission-এর cumulative energy। এটি একটি UAV-এর single-battery limit নয় এবং একাধিক mission-এর কারণে ১০০০ J-এর বেশি হতে পারে।

\begin{table}[htbp]
\centering
\caption{V1 energy summary এবং energy efficiency।}
\label{tab:energy-summary-bn}
\scriptsize
\rowcolors{2}{lightgray}{white}
\begin{tabular}{@{}lrrrr@{}}
\toprule
Environment & Total energy/run (J) & Energy/task (J) & Energy/parcel (J) & Distance/run (m) \\
\midrule
DU outdoor & 6142.4 $\pm$ 2417.8 & 310.2 & 199.3 & 5928.2 \\
Synthetic & 5319.7 $\pm$ 2101.8 & 134.4 & 84.3 & 5041.4 \\
\bottomrule
\end{tabular}
\end{table}

এই summary-তে প্রতিটি environment-এর ১২০টি run-এর average ব্যবহার করা হয়েছে। Total energy-এর পাশে $\pm$ value standard deviation; energy/task এবং energy/parcel হলো run-level ratio-এর average। DU outdoor-এ Synthetic-এর তুলনায় cumulative energy প্রায় ১৫.৫ শতাংশ বেশি এবং delivered parcel প্রতি energy প্রায় ২.৪ গুণ বেশি। এর প্রধান কারণ DU outdoor-এ route distance ও route geometry কঠিন, এবং কম parcel deliver হওয়ায় একই ধরনের mission energy কম completed service-এর মধ্যে ভাগ হয়।

Task load ৩০ থেকে ৯০ করলে average total energy Synthetic-এ ৪২১৫.৭ J থেকে ৬১১৫.১ J এবং DU outdoor-এ ৫৫১০.৪ J থেকে ৬৫২২.৮ J হয়েছে। একইভাবে UAV সংখ্যা ৩ থেকে ১০ করলে cumulative system energy Synthetic-এ ২৮১৯.১ J থেকে ৭৪৯৬.০ J এবং DU outdoor-এ ৩০৪৮.৭ J থেকে ৯০৫৬.৯ J হয়েছে। এই বৃদ্ধি বেশি mission এবং বেশি completed service-এর cumulative effect; এটিকে একটি UAV-এর battery depletion হিসেবে interpret করা যাবে না।

\section{Simulation-এর ফলাফল}

\subsection{Baseline: ৩০ task এবং ৩ UAV}

৩০ task এবং ৩ UAV configuration-কে সহজ reference point হিসেবে নেওয়া হয়েছে। নিচের value-গুলো ১০টি seed-এর average; $\pm$-এর পরের value standard deviation।

\begin{table}[htbp]
\centering
\caption{৩০ task এবং ৩ UAV baseline result।}
\label{tab:baseline-bn}
\scriptsize
\rowcolors{2}{lightgray}{white}
\begin{tabular}{@{}lrrrrrrrrr@{}}
\toprule
Environment & Objective & Served & Task rate & Req. parcel & Dropped & Parcel rate & Deferred & Delay & Runtime \\
\midrule
DU outdoor & 5.107 $\pm$ 1.501 & 10.500 $\pm$ 3.100 & 35.000\% & 48.8 & 14.9 & 30.533\% & 19.500 $\pm$ 3.100 & 34.333 & 11.485 \\
Synthetic & 11.981 $\pm$ 1.484 & 22.600 $\pm$ 1.776 & 75.333\% & 48.6 & 34.5 & 70.988\% & 7.400 $\pm$ 1.776 & 23.656 & 2.069 \\
\bottomrule
\end{tabular}
\end{table}

এই baseline-এ Synthetic-এ গড়ে ২২.৬টি task এবং ৩৪.৫টি parcel complete হয়েছে, যেখানে DU outdoor-এ হয়েছে ১০.৫টি task এবং ১৪.৯টি parcel। ফলে Synthetic-এর task rate ও parcel delivery rate যথাক্রমে ৭৫.৩৩% ও ৭০.৯৯%, কিন্তু DU outdoor-এ তা ৩৫.০০% ও ৩০.৫৩%। একই সঙ্গে DU outdoor-এ delay এবং runtime বেশি—বর্তমান geometry ও route search-এর কারণে এই case-টি বেশি কঠিন।

\subsection{Task load বাড়ালে কী হয়}

প্রতিটি row একই environment এবং task count-এর জন্য ৩, ৫, ৮ ও ১০ UAV configuration-এর ফলাফল একসাথে average করে তৈরি করা হয়েছে; প্রতিটি configuration-এ ১০টি seed ব্যবহার করা হয়েছে। তাই এই table-এ কোনো একটি নির্দিষ্ট UAV সংখ্যা নয়, চারটি fleet size-এর combined effect দেখানো হয়েছে।

\begin{table}[htbp]
\centering
\caption{Task load পরিবর্তনের প্রভাব।}
\label{tab:task-effect-bn}
\scriptsize
\rowcolors{2}{lightgray}{white}
\begin{tabular}{@{}llrrrrrrrr@{}}
\toprule
Environment & Tasks & Objective & Served & Task rate & Dropped & Parcel rate & Peak queue & Planner (s) & Route \\
\midrule
DU outdoor & 30 & 8.260 $\pm$ 2.610 & 16.950 $\pm$ 5.198 & 56.500\% & 25.5 & 52.254\% & 13.3 & 13.639 & 96.364\% \\
DU outdoor & 60 & 11.318 $\pm$ 4.378 & 22.025 $\pm$ 9.003 & 36.708\% & 35.1 & 34.674\% & 38.0 & 21.273 & 96.981\% \\
DU outdoor & 90 & 11.726 $\pm$ 4.584 & 21.400 $\pm$ 8.311 & 23.778\% & 34.0 & 22.487\% & 68.6 & 25.757 & 96.774\% \\
Synthetic & 30 & 14.995 $\pm$ 2.161 & 26.200 $\pm$ 2.564 & 87.333\% & 41.7 & 85.751\% & 5.3 & 1.340 & 97.678\% \\
Synthetic & 60 & 25.013 $\pm$ 6.752 & 43.475 $\pm$ 12.471 & 72.458\% & 72.8 & 69.291\% & 19.2 & 3.800 & 95.829\% \\
Synthetic & 90 & 30.430 $\pm$ 9.562 & 51.525 $\pm$ 18.664 & 57.250\% & 81.6 & 52.190\% & 39.9 & 6.053 & 94.701\% \\
\bottomrule
\end{tabular}
\end{table}

Task load বাড়ার সঙ্গে সঙ্গে peak queue দ্রুত বেড়েছে এবং task ও parcel service rate কমেছে। DU outdoor-এ ৩০ task থেকে ৯০ task-এ task service rate ৫৬.৫০ শতাংশ থেকে ২৩.৭৮ শতাংশে এবং parcel delivery rate ৫২.২৫ শতাংশ থেকে ২২.৪৯ শতাংশে নেমেছে। Synthetic-এ একই পরিবর্তনে task service rate ৮৭.৩৩ শতাংশ থেকে ৫৭.২৫ শতাংশে এবং parcel delivery rate ৮৫.৭৫ শতাংশ থেকে ৫২.১৯ শতাংশে নেমেছে। Objective value বাড়লেও তা বেশি task load-এর কারণে মোট accumulated value; এটি service quality বাড়ার সরাসরি প্রমাণ নয়।

এখানে dropped parcel-এর average সবসময় served task-এর সঙ্গে একই অনুপাতে বাড়েনি, কারণ প্রতিটি task-এ parcel সংখ্যা এক নয়। তাই task rate-এর পাশাপাশি parcel delivery rate report করা জরুরি। মোট task বাড়ায় deferred task এবং parcel backlog-এর উপর চাপ স্পষ্টভাবে বেড়েছে।

\subsection{UAV সংখ্যা বাড়ালে কী হয়}

প্রতিটি row একই environment এবং UAV count-এর জন্য ৩০, ৬০ ও ৯০ task configuration-এর ফলাফল একসাথে average করে তৈরি করা হয়েছে; প্রতিটি task load-এ ১০টি seed ব্যবহার করা হয়েছে। অর্থাৎ এখানে UAV সংখ্যার effect দেখানো হচ্ছে, task load fixed রাখা হয়নি।

\begin{table}[htbp]
\centering
\caption{UAV সংখ্যা পরিবর্তনের প্রভাব।}
\label{tab:fleet-effect-bn}
\scriptsize
\rowcolors{2}{lightgray}{white}
\begin{tabular}{@{}llrrrrrrr@{}}
\toprule
Environment & UAV & Objective & Served & Task rate & Dropped & Parcel rate & Deferred & Runtime (s) \\
\midrule
DU outdoor & 3 & 5.770 $\pm$ 1.390 & 10.833 $\pm$ 2.183 & 18.056\% & 16.2 & 16.097\% & 49.2 & 13.331 \\
DU outdoor & 5 & 8.922 $\pm$ 1.934 & 17.100 $\pm$ 2.820 & 28.500\% & 26.3 & 26.153\% & 42.9 & 17.240 \\
DU outdoor & 8 & 12.530 $\pm$ 2.740 & 24.233 $\pm$ 4.289 & 40.389\% & 38.3 & 38.168\% & 35.8 & 23.273 \\
DU outdoor & 10 & 14.517 $\pm$ 3.641 & 28.333 $\pm$ 6.504 & 47.222\% & 45.4 & 45.204\% & 31.7 & 27.058 \\
Synthetic & 3 & 15.144 $\pm$ 2.919 & 25.233 $\pm$ 3.025 & 42.056\% & 37.6 & 36.408\% & 34.8 & 2.726 \\
Synthetic & 5 & 21.416 $\pm$ 4.964 & 36.300 $\pm$ 7.293 & 60.500\% & 56.9 & 55.047\% & 23.7 & 3.453 \\
Synthetic & 8 & 27.414 $\pm$ 8.774 & 47.700 $\pm$ 15.434 & 79.500\% & 78.9 & 76.298\% & 12.3 & 4.351 \\
Synthetic & 10 & 29.941 $\pm$ 10.681 & 52.367 $\pm$ 19.728 & 87.278\% & 88.1 & 85.231\% & 7.6 & 4.403 \\
\bottomrule
\end{tabular}
\end{table}

UAV সংখ্যা ৩ থেকে ১০ করলে দুই environment-এই task service এবং parcel delivery rate বাড়ে, deferred queue কমে। DU outdoor-এ parcel delivery rate ১৬.১০ শতাংশ থেকে ৪৫.২০ শতাংশে এবং Synthetic-এ ৩৬.৪১ শতাংশ থেকে ৮৫.২৩ শতাংশে বেড়েছে। তবে computation time ও candidate search চাপও বেড়েছে, কারণ বেশি UAV থাকলে assignment এবং task order-এর সম্ভাব্য combination সংখ্যাও বেড়ে যায়। তাই UAV বাড়ানো service-এর জন্য উপকারী হলেও এটি computationally free নয়।

\clearpage
\begin{landscape}
\subsection{সম্পূর্ণ task load এবং fleet matrix}

নিচের table-এ প্রতিটি environment, task load এবং UAV configuration-এর দশটি seed-এর average দেওয়া হয়েছে। এটি experiment-এর সম্পূর্ণ performance landscape।

\scriptsize
\begin{longtable}{@{}llrrrrrrrrrrr@{}}
\caption{সম্পূর্ণ V1 result matrix।}\label{tab:full-matrix-bn}\\
\toprule
Environment & Task & UAV & Objective & Served & Task rate & Req. parcel & Dropped & Parcel rate & Deferred & Delay & Runtime & Route \\
\midrule
\endfirsthead
\toprule
Environment & Task & UAV & Objective & Served & Task rate & Req. parcel & Dropped & Parcel rate & Deferred & Delay & Runtime & Route \\
\midrule
\endhead
DU outdoor & 30 & 3 & 5.107 $\pm$ 1.501 & 10.500 $\pm$ 3.100 & 35.000\% & 48.8 & 14.9 & 30.533\% & 19.500 $\pm$ 3.100 & 34.333 & 11.485 & 99.266\% \\
DU outdoor & 30 & 5 & 7.472 $\pm$ 1.423 & 15.900 $\pm$ 3.315 & 53.000\% & 48.8 & 23.3 & 47.746\% & 14.100 $\pm$ 3.315 & 33.314 & 13.509 & 97.982\% \\
DU outdoor & 30 & 8 & 9.873 $\pm$ 1.565 & 20.200 $\pm$ 3.120 & 67.333\% & 48.8 & 31.0 & 63.525\% & 9.800 $\pm$ 3.120 & 28.357 & 15.143 & 95.013\% \\
DU outdoor & 30 & 10 & 10.589 $\pm$ 1.461 & 21.200 $\pm$ 2.741 & 70.667\% & 48.8 & 32.8 & 67.213\% & 8.800 $\pm$ 2.741 & 25.973 & 14.422 & 93.954\% \\
DU outdoor & 60 & 3 & 5.946 $\pm$ 1.272 & 10.800 $\pm$ 1.619 & 18.000\% & 101.3 & 16.7 & 16.486\% & 49.200 $\pm$ 1.619 & 33.638 & 12.684 & 99.316\% \\
DU outdoor & 60 & 5 & 9.490 $\pm$ 1.524 & 17.800 $\pm$ 1.989 & 29.667\% & 101.3 & 27.9 & 27.542\% & 42.200 $\pm$ 1.989 & 33.207 & 18.059 & 97.141\% \\
DU outdoor & 60 & 8 & 13.633 $\pm$ 1.996 & 26.900 $\pm$ 3.035 & 44.833\% & 101.3 & 42.9 & 42.349\% & 33.100 $\pm$ 3.035 & 32.947 & 23.971 & 96.671\% \\
DU outdoor & 60 & 10 & 16.203 $\pm$ 2.625 & 32.600 $\pm$ 4.993 & 54.333\% & 101.3 & 53.0 & 52.320\% & 27.400 $\pm$ 4.993 & 32.623 & 30.386 & 96.158\% \\
DU outdoor & 90 & 3 & 6.258 $\pm$ 1.251 & 11.200 $\pm$ 1.687 & 12.444\% & 151.2 & 16.9 & 11.177\% & 78.800 $\pm$ 1.687 & 32.996 & 15.823 & 99.210\% \\
DU outdoor & 90 & 5 & 9.805 $\pm$ 2.035 & 17.600 $\pm$ 2.875 & 19.556\% & 151.2 & 27.6 & 18.254\% & 72.400 $\pm$ 2.875 & 31.723 & 20.151 & 97.796\% \\
DU outdoor & 90 & 8 & 14.084 $\pm$ 2.424 & 25.600 $\pm$ 3.502 & 28.444\% & 151.2 & 41.1 & 27.183\% & 64.400 $\pm$ 3.502 & 30.099 & 30.704 & 96.515\% \\
DU outdoor & 90 & 10 & 16.758 $\pm$ 2.796 & 31.200 $\pm$ 4.237 & 34.667\% & 151.2 & 50.4 & 33.333\% & 58.800 $\pm$ 4.237 & 30.517 & 36.365 & 95.366\% \\
Synthetic & 30 & 3 & 11.981 $\pm$ 1.484 & 22.600 $\pm$ 1.776 & 75.333\% & 48.6 & 34.5 & 70.988\% & 7.400 $\pm$ 1.776 & 23.656 & 2.069 & 96.483\% \\
Synthetic & 30 & 5 & 15.303 $\pm$ 1.102 & 27.200 $\pm$ 1.398 & 90.667\% & 48.6 & 43.8 & 90.123\% & 2.800 $\pm$ 1.398 & 16.152 & 1.118 & 99.377\% \\
Synthetic & 30 & 8 & 16.259 $\pm$ 1.138 & 27.500 $\pm$ 1.434 & 91.667\% & 48.6 & 44.2 & 90.947\% & 2.500 $\pm$ 1.434 & 11.406 & 1.102 & 98.288\% \\
Synthetic & 30 & 10 & 16.436 $\pm$ 1.117 & 27.500 $\pm$ 1.434 & 91.667\% & 48.6 & 44.2 & 90.947\% & 2.500 $\pm$ 1.434 & 10.370 & 1.078 & 97.587\% \\
Synthetic & 60 & 3 & 15.646 $\pm$ 1.861 & 25.400 $\pm$ 2.459 & 42.333\% & 105.1 & 39.1 & 37.203\% & 34.600 $\pm$ 2.459 & 21.600 & 2.862 & 99.258\% \\
Synthetic & 60 & 5 & 22.837 $\pm$ 2.172 & 39.900 $\pm$ 3.843 & 66.500\% & 105.1 & 63.9 & 60.799\% & 20.100 $\pm$ 3.843 & 21.642 & 4.039 & 96.426\% \\
Synthetic & 60 & 8 & 29.704 $\pm$ 2.663 & 53.300 $\pm$ 3.199 & 88.833\% & 105.1 & 92.0 & 87.536\% & 6.700 $\pm$ 3.199 & 18.928 & 4.663 & 94.019\% \\
Synthetic & 60 & 10 & 31.863 $\pm$ 1.716 & 55.300 $\pm$ 1.767 & 92.167\% & 105.1 & 96.3 & 91.627\% & 4.700 $\pm$ 1.767 & 15.013 & 3.644 & 94.788\% \\
Synthetic & 90 & 3 & 17.805 $\pm$ 1.590 & 27.700 $\pm$ 2.406 & 30.778\% & 156.4 & 39.3 & 25.128\% & 62.300 $\pm$ 2.406 & 20.887 & 3.246 & 99.802\% \\
Synthetic & 90 & 5 & 26.108 $\pm$ 2.291 & 41.800 $\pm$ 3.824 & 46.444\% & 156.4 & 63.0 & 40.281\% & 48.200 $\pm$ 3.824 & 20.039 & 5.201 & 96.675\% \\
Synthetic & 90 & 8 & 36.280 $\pm$ 2.869 & 62.300 $\pm$ 5.498 & 69.222\% & 156.4 & 100.4 & 64.194\% & 27.700 $\pm$ 5.498 & 19.906 & 7.286 & 93.854\% \\
Synthetic & 90 & 10 & 41.525 $\pm$ 2.748 & 74.300 $\pm$ 4.218 & 82.556\% & 156.4 & 123.8 & 79.156\% & 15.700 $\pm$ 4.218 & 19.960 & 8.488 & 92.269\% \\
\bottomrule
\end{longtable}

\textbf{এই matrix-এর key finding:} Synthetic ৩০-task case-এ ৫ বা তার বেশি UAV ব্যবহার করলে parcel delivery rate প্রায় ৯০\% হয়, কিন্তু DU outdoor-এ ৯০ task এবং ১০ UAV থাকলেও তা ৩৩.৩৩\%-এ সীমিত থাকে। UAV বাড়ালে service বাড়ে, তবে বিশেষ করে DU outdoor-এ runtime-ও বাড়ে; তাই workload, fleet size এবং environment একসাথে বিবেচনা করা জরুরি।
\end{landscape}

\section{Runtime এবং প্রধান bottleneck}

V1-এর সবচেয়ে গুরুত্বপূর্ণ computational finding হলো planner-এর সময়ের প্রায় পুরো অংশ route evaluation-এ ব্যয় হয়েছে।

\begin{table}[htbp]
\centering
\caption{Environment অনুযায়ী runtime এবং search evidence।}
\label{tab:runtime-bn}
\scriptsize
\rowcolors{2}{lightgray}{white}
\begin{tabular}{@{}lrrrrrr@{}}
\toprule
Environment & Mean planner s & Min s & Max s & Mean candidate & Max candidate & Route share \\
\midrule
DU outdoor & 20.223 & 3.325 & 62.198 & 249476.700 & 50000 & 96.754\% \\
Synthetic & 3.731 & 0.157 & 14.430 & 65210.608 & 50000 & 95.440\% \\
\bottomrule
\end{tabular}
\end{table}

সম্পূর্ণ experiment-এ ৬১২টি dispatch ৫০,০০০ candidate search limit-এ পৌঁছেছে। এর অর্থ হলো ওই dispatch-গুলোতে V1 নির্ধারিত bound-এর মধ্যে search থামিয়েছে; এগুলোকে global optimum হিসেবে দাবি করা যাবে না। বিশেষ করে ১০ UAV এবং বড় task load case-এ এই truncation বেশি হয়েছে।

এই evidence থেকে V2-এর জন্য সবচেয়ে গুরুত্বপূর্ণ improvement target হলো:

\begin{itemize}
\item V1-এর basic route cache-এর hit rate এবং reuse record করা, এবং cache key ও lifetime আরও নির্ভরযোগ্য করা;
\item infeasible candidate শুরুতেই বাদ দেওয়া;
\item task এবং UAV selection-এর search space ছোট করা;
\item fixed active route-এর information পরবর্তী dispatch-এ reuse করা;
\item bounded brute-force-এর পরিবর্তে scalable candidate selection ব্যবহার করা।
\end{itemize}

\section{V1 result-এর অর্থ এবং limitation}

এই experiment থেকে প্রধান findings হলো:

\begin{enumerate}
\item আগের ছয় task-এর ছোট experiment queue pressure বোঝার জন্য যথেষ্ট ছিল না। ৩০, ৬০ এবং ৯০ task এবং ৩, ৫, ৮ ও ১০ UAV-এর heavy matrix-এ queue এবং deferred work স্পষ্টভাবে দেখা গেছে।
\item High, medium এবং low phase-এর ৩:২:১ structure প্রতিটি run-এ সঠিকভাবে reproduce হয়েছে।
\item Synthetic এবং DU outdoor result-এর পার্থক্য environment provenance-এর সঙ্গে পড়তে হবে: DU case-এ OSM-derived layout-এর উপর derived damage geometry ব্যবহৃত হয়েছে, actual post-earthquake field reconstruction নয়।
\item Task demand individual request হিসেবে তৈরি হয়েছে এবং item weight, parcel count, damaged-footprint location ও accessible drop-off waypoint-এর নিয়ম result-কে প্রভাবিত করেছে। তাই task এবং parcel rate আলাদা করে দেখা প্রয়োজন।
\item UAV সংখ্যা বাড়ালে objective value, task service rate এবং parcel delivery rate বাড়ে এবং queue কমে, কিন্তু planner-এর search space এবং runtime বাড়ে।
\item Task count একা delivered service বোঝানোর জন্য যথেষ্ট নয়। একটি task-এ একাধিক parcel থাকতে পারে, তাই dropped parcel count এবং parcel delivery rate আলাদা করে report করা হয়েছে।
\item প্রতি parcel-এর জন্য fixed ৩০ second drop-off time V1-এর একটি স্পষ্ট timing এবং visualization assumption। এটি field-calibrated unloading measurement নয়।
\item DU outdoor environment এই implementation-এ Synthetic-এর তুলনায় ধীর এবং কম serviceable হয়েছে। এটিকে universal real-world law হিসেবে interpret করা যাবে না; এটি বর্তমান geometry এবং routing model-এর measured ফলাফল।
\item Route evaluation V1-এর প্রধান runtime bottleneck।
\item সব ২৪০টি run-এ planned safe-return rate ১০০ শতাংশ হয়েছে। এটি planner-এর return এবং reserve-energy feasibility filter-এর ফল; real-world flight reliability-এর estimate নয়।
\item V1-এর basic route cache ইতিমধ্যে repeated route-leg construction কমাতে ব্যবহৃত হয়েছে; V2-তে cache reuse এবং hit-rate আলাদাভাবে evaluate করা উচিত।
\item Energy V1-এ route feasibility এবং reserve check-এ ব্যবহৃত হয়েছে, কিন্তু current energy coefficients field-calibrated নয় এবং event-by-event battery trajectory আলাদাভাবে record করা হয়নি।
\item V1 service এবং computational behaviour বোঝার জন্য একটি useful reference baseline, কিন্তু এটি final mathematical optimizer নয়।
\end{enumerate}

V1-এর সীমাবদ্ধতা:

\begin{itemize}
\item Assignment এবং ordering-এর জন্য bounded brute-force ব্যবহার করা হয়েছে।
\item Search limit-এ পৌঁছানো result global optimum নয়।
\item Recourse এবং rerouting V1-এ নেই।
\item ২ kg payload, প্রতি parcel ৩০ second drop-off time, ৫ মিনিট turnaround এবং sufficient base supply বর্তমান simulation assumption।
\item Task-এর প্রতিটি item-এর quantity V1-এ ০ অথবা ১; formulation-এর মতো arbitrary integer quantity বা time-varying remaining demand নেই।
\item Payload utilization metric fleet parcel-count ratio হিসেবে সংজ্ঞায়িত; এটি per-UAV payload mass utilization নয়।
\item Energy coefficient এবং operational parameters এখনো field-calibrated নয়।
\item Total energy cumulative mission energy হিসেবে record করা হয়েছে; per-UAV battery state এবং recharge curve V1-এ explicitly model করা হয়নি।
\item Service rate-কে বাস্তব field reliability estimate হিসেবে ব্যবহার করা যাবে না।
\end{itemize}

\section{উপসংহার এবং পরবর্তী কাজ}

V1 simulation একটি complete two-hour, dynamic, multi-seed এবং multi-scenario experiment হিসেবে কাজ করেছে। এটি শুধু assignment এবং route তৈরি করেনি, বরং task arrival, queue, UAV lifecycle, parcel drop-off, return এবং runtime bottleneck-ও record করেছে।

ফলাফল দেখিয়েছে যে বর্তমান V1 low এবং moderate workload-এ কাজ করতে পারে, কিন্তু task load বাড়লে queue দ্রুত বৃদ্ধি পায় এবং task ও parcel delivery rate কমে যায়। UAV বাড়ালে objective value এবং delivery performance উন্নত হয়, কিন্তু brute-force search-এর computational cost-ও বাড়ে।

তাই পরবর্তী ধাপে প্রথমে একটি efficient V2 planner তৈরি করা উচিত। V2-তে
existing route cache-এর reuse আরও শক্তিশালী করা, early pruning এবং scalable
candidate search যোগ করে V1-এর সঙ্গে সরাসরি comparison করতে হবে। তবে recourse
experiment শুরু করার আগে UAV state model শক্তিশালী করা প্রয়োজন। নতুন task এলে
active UAV-এর current position, remaining route, onboard item এবং remaining
battery জানা না থাকলে formulation-এর replacement feasibility যাচাই করা যাবে না।
তাই stateful battery-and-inventory update এবং true task-replacement recourse
V2-এর প্রধান research extension হবে। Efficiency improvement এবং recourse
result আলাদা করে report করা হবে।

\section{Reproducibility artifacts}

এই report-এর ফলাফল নিচের local artifacts থেকে reproduce করা যায়:

{\scriptsize
\begin{itemize}[itemsep=1pt,topsep=2pt]
\item \code{Simulation/results/pgbm_v1_two_hour_experiments/raw_metrics/pgbm_v1_two_hour_matrix_metrics_tasks_30_60_90_uavs_3_5_8_10_seeds_101_110.csv}: ২৪০টি raw result row, objective এবং parcel metrics সহ;
\item \code{Simulation/results/pgbm_v1_two_hour_experiments/manifests/pgbm_v1_two_hour_matrix_manifest.json}: experiment configuration এবং coverage;
\item \code{Simulation/run_pgbm_v1_experiment_matrix.py}: heavy matrix run করার script;
\item \code{Simulation/pgbm_sim/experiment.py}: metrics row তৈরি করার logic;
\item \code{Simulation/pgbm_sim/event_simulation.py}: ১২০ মিনিটের event clock এবং queue execution।
\end{itemize}
}

\end{document}
